% SGA 3, 第 VI_B 篇, 第 1 轮英文重构.
% 范围: 推论 5.6.2 全文, 含底本编者注 43.
% 底本: Polo--Gille Exp6B-13oct24.pdf, 分篇第 32--33 页
% (合订本第 373--374 页; 原版页边印刷页码 355, 重版本页眉 362--363).
% 引理 5.6.3 始于分篇第 33 页 / 合订本第 374 页.
% OCR 与 Reinhold 英文 Markdown 仅为比对材料.
% 本单元无图.

\medskip
\phantomsection\label{cor:sga3-VIB-5.6.2}%
\noindent\textbf{推论 5.6.2.---}
\emph{设 \(S\), \(G\) 和 \(X\) 如
\hyperref[thm:sga3-VIB-5.3]{5.3} 的预备假设中所述,
\(U\) 是 \(X\) 的开集, \(s\in S\), \(A\) 是 \(X_s\) 的有限子集.
假设 \(U_s\) 在 \(X_s\) 中稠密,
且 \(X_s\) 在 \(\kappa(s)\) 上局部有限型.}%
\footnote{\label{note:sga3-VIB-43}编者注: 我们增加了关于 \(X_s\) 的假设,
并考虑到新增的 \SGARefLink{lem:sga3-VIB-5.6.2.0}{5.6.2.0},
简化并改正了证明. 此外, 证明表明, 仅假设 \(U_s\)
在 \(X_s\) 的每个不可约分支中稠密时, 结论仍成立.}

\emph{则存在态射 \(f:S''\to S\) 和 \(h:S''\to G\),
其中 \(f\) 分解为有限满射态射 \(S''\to S'\)
后接平展态射 \(S'\to S\), 使得 \(A\) 在 \(X\times_S S''\)
中的逆像 \(A''\) \textup{(即在
\(X_s\times_{\operatorname{Spec}\kappa(s)}S''_s\) 中的逆像)} 包含在}
\[
  \ell_h^{-1}(U_{S''}),
\]
\emph{中, 其中 \(\ell_h\) 表示由元素
\(h\in G(S'')\) 定义的 \(X_{S''}=X\times_S S''\) 的平移.}

\emph{证明.}
由于 \(X_s\) 在 \(\kappa(s)\) 上局部有限型,
\(X_s\) 的连通分支开且不可约
(参见 VI\(_A\), \SGARefLink{prop:sga3-VIA-2.5.4}{2.5.4}).
所以 \(U_s\) 在 \(X_s\) 的每个不可约分支中稠密.

因此由\hyperref[lem:sga3-VIB-5.6.2.0]{引理 5.6.2.0}, 存在闭点
\[
  g\in G_s^0
\]
使得对每个位于 \(A\) 的某点上的点 \(a'\in X_{\kappa(g)}\), 有
\[
  g\cdot a'\in U_s\otimes_{\kappa(s)}\kappa(g)
\]

由\hyperref[lem:sga3-VIB-5.6.1]{引理 5.6.1},
存在态射 \(f:S''\to S\) 和 \(h:S''\to G\),
其中 \(f\) 分解为有限满射态射 \(S''\to S'\)
后接平展态射 \(S'\to S\),
使 \(f^{-1}(s)\) 仅由一点 \(s_0\) 组成, 且
\[
  h(s_0)=g,\qquad \kappa(s_0)=\kappa(g).
\]
由于 \(S''_s\) 有限且只有一点 \(s_0\), 典范闭浸入
\[
  X_s\otimes_{\kappa(s)}\kappa(s_0)
  \longrightarrow
  X_s\times_{\operatorname{Spec}\kappa(s)}S''_s
\]
是幂零的, 因而在底空间之间诱导同胚.
\footnote{Canon 底本注 (SGA-SRC-000863): 重版本将这些概形等同.
有限的单点纤维未必既约; 从其剩余域出发的典范映射实为幂零闭浸入,
这恰好足以支持此处的点集论证.}
用 \(A''\) 表示 \(A\) 在右端底空间中的逆像.
通过此同胚及 \(\kappa(s_0)=\kappa(g)\),
它与 \(A\) 在
\[
  X_s\otimes_{\kappa(s)}\kappa(s_0)
  =X_s\otimes_{\kappa(s)}\kappa(g).
\]
中的逆像等同. 于是 \(X_{S''}\) 的平移 \(\ell_h\) 将 \(A''\) 映入
\[
  U_s\otimes_{\kappa(s)}\kappa(s_0).
\]
